\section{Preliminaries} \label{sec:prelims}

\myparagraph{Basic notions} A \emph{graph} will always mean a finite directed graph without self-loops.
A \emph{colored graph} is a graph where nodes have a unique node label, and every pair of distinct nodes is connected by a unique edge (in one direction or the other), with node and edge vocabularies distinct.
A \emph{multigraph} is a finite directed graph without self-loops that allows multiple edges (including no edge) between a pair of distinct nodes.
A \emph{colored multigraph} is defined similarly, but requiring each node to have a unique label.
The multiplicity of a multigraph is defined as the maximum edge multiplicity within it.

The $i^{th}$ unit vector, denoted by $\unit{i}$,
has $i^{th}$ entry $1$, and other entries are $0$.
We define $\bfzero$ as the vector where all entries are $0$,
and $\bfone$ as the vector that all entries are $1$.
For  $\bfv, \bfu \in \bbN^n$,
$\bfv \vle \bfu$ if, for $1 \le i \le n$, $\bfv_i \le \bfu_i$
Let $\projp{n}{\bfv}$ denote the projection operator that maps $\bfv$ to its first $n$ entries.
Let $\norm{\bfA}$ and $\norm{\bfv}$ denote the maximal absolute values of the entries in $\bfA$ and $\bfv$, respectively.

\myparagraph{Linear algebra and Kleene star}
An \emph{integer linear constraint} is an inequality $\sum_{i} a_i \cdot x_i \leq c$
involving integer coefficients $a_i$, non-negative integer variables $x_i$, and an integer constant $c$.
The set of solutions of such a constraint is called a \emph{modulus-free semilinear set}, while
if we include equations $x = c \mod p$ for integers $c, p>0$ it is a \emph{semilinear set}.
A constraint is \emph{homogeneous} if $c = 0$.

For an integer linear system $\cQ(\bfx): \bfA \bfx = \bfc$,
we denote by $\sol{\cQ(\bfx)}$ the set of solutions in $\bbN$,
i.e., $\sol{\cQ(\bfx)} := \setc{\bfv \in \bbN^n}{\bfA\bfv = \bfc}$, and similarly for a Boolean combination of linear systems.
Note that if $\cQ(\bfx)$ is feasible,
then it has a solution in which the number of non-zero assignments is bounded, and the values of the assignments are also bounded.

\begin{lemma}[\cite{papa}, \cite{caratheodory-integer}, also see Corollary 2.2 in~\cite{twovarpres}]\label{lem:ilp}
    There are constants $c_1, c_2 \in \bbN$ such that for every integer linear system $\cQ(\bfx)$,
    if $\cQ(\bfx)$ admits a solution in $\bbN$, then it admits a solution in $\bbN$ in which
    the number of variables assigned non-zero values is at most $c_1t\log\left(c_2tM\right)$ and
    every variable is assigned a value bounded by $c_1t(tM)^{c_2t}$,
    where $t$ is the number of constraints in $\cQ(\bfx)$ and $M$ is the maximal absolute value of the coefficients in $\cQ(\bfx)$.
\end{lemma}

